\documentclass{article}
\usepackage{PRIMEarxiv} % Core package for the PrimeAI Template
\usepackage{amsmath, amssymb, amsthm, bm, dcolumn} % Add amsthm here for the proof environment
\usepackage[numbers,sort&compress]{natbib} % Natbib for citations
\usepackage{graphicx} % For high-quality images
\usepackage{multicol} % Multi-column figures/tables
\usepackage{hyperref} % Hyperlinks
\usepackage{listings} % Code listings
\usepackage{authblk} % For structured affiliations
% Define theorem style (optional)
\theoremstyle{plain}
\newtheorem{theorem}{Theorem}
\renewcommand{\qedsymbol}{}

% Custom commands for primes and qubit encoding
\newcommand{\primeQubit}[1]{|p_{#1}\rangle}
\newcommand{\primeGate}[1]{U_{p_{#1}}}

\usepackage{wrapfig}
\usepackage[pscoord]{eso-pic}
\usepackage[fulladjust]{marginnote}
\reversemarginpar

% Typesetting improvements without footnote patching
\usepackage[protrusion=true, expansion=true, tracking=false]{microtype}
\microtypecontext{spacing=nonfrench}

% Line numbers
\usepackage[right]{lineno}

% Text layout - adjust as needed
\raggedright
\setlength{\parindent}{0.5cm}
\textwidth 5.25in 
\textheight 8.75in

% Set double spacing
\usepackage{setspace} 
\doublespacing

% Adjust width for specific content
\usepackage{changepage}

% Adjust caption style
\usepackage[aboveskip=1pt,labelfont=bf,labelsep=period,singlelinecheck=off]{caption}

% Remove brackets from references
\makeatletter
\renewcommand{\@biblabel}[1]{\quad#1.}
\makeatother

% Header, footer, and page numbers
\usepackage{lastpage,fancyhdr}
\pagestyle{fancy}
\fancyhf{}
\fancyhead[L]{Preprint - PrimeAI Enhanced Template}
\fancyfoot[C]{\scriptsize Multiplicity Theory © 2024 - Citizen Gardens \\ Licensed Under MIT and CC BY-NC-SA 4.0.}
\fancyfoot[R]{Page \thepage\ of \pageref{LastPage}}
\renewcommand{\footrule}{\hrule height 2pt \vspace{2mm}}



\begin{document}

\title{Integrating Pythagorean Triplets and Fibonacci Sequences, with Multiplicity Theory}
\author{Miroslav Zidek, Ryan O. Van Gelder}
\date{\today}
\maketitle
\begin{abstract}
This paper presents an enhanced framework that integrates Pythagorean triplets, Fibonacci sequences, and Multiplicity Theory. The framework is built on dynamic multiplicity equations, prime-based encoding, recursive feedback, and tensor dynamics. It is extended to practical applications in cryptography, astrophysics, and educational tools, while maintaining mathematical consistency and interdisciplinary applicability.
\end{abstract}

\section{Core Dynamic Multiplicity Equation}
The dynamic multiplicity equation, incorporating nonlinearity, memory, and geometric feedback, is given as:
\begin{equation}
    \frac{\partial \rho_k}{\partial t} = \alpha_k(t) \rho_k + \beta_k(t) I_k + \gamma_k(t) \sum_{j} T_{kj} \rho_j + \lambda(t) (\Omega_B(\rho) + \Omega_{FS}(\rho)) + \eta_k \rho_k^2 + \zeta_k \sum_{m,n} C_{mn} \rho_m \rho_n + \mu_k M_k(t),
\end{equation}
where:
\begin{itemize}
    \item $\alpha_k(t), \beta_k(t), \gamma_k(t), \lambda(t)$: Time-dependent coefficients.
    \item $T_{kj}$: Tensor coupling terms.
    \item $\Omega_B(\rho), \Omega_{FS}(\rho)$: Geometric feedback terms.
    \item $\eta_k \rho_k^2$: Nonlinear self-interaction term.
    \item $C_{mn}$: Correlation tensors for multi-scale interactions.
    \item $\mu_k M_k(t)$: Long-term memory effects.
\end{itemize}

\section{Prime-Based Encoding}
Define a prime encoding function $P(x, t)$ for Fibonacci numbers and Pythagorean triplets:
\begin{equation}
    P(x, t) = \prod_{i=1}^{N} p_i^{f_i(t)},
\end{equation}
where:
\begin{itemize}
    \item $p_i$: Prime numbers.
    \item $f_i(t) = F_n \mod m$: Fibonacci coefficients modulated by a recursive function.
\end{itemize}
The dynamic encoding evolves as:
\begin{equation}
    P(x, t) = P_{\text{base}}(x) \cdot F(t),
\end{equation}
with stochastic noise $\epsilon(t)$:
\begin{equation}
    F(t) = 1 + \epsilon(t).
\end{equation}

\section{Tensor and Hypergraph Dynamics}
Pythagorean triplets and Fibonacci sequences are modeled as hypergraph nodes with adjacency tensor:
\begin{equation}
    T_{ijk} = \prod_{p \in \text{Hyperedge}(i,j,k)} p_{F_i + F_j + F_k}.
\end{equation}
Eigenvalue dynamics track system stability:
\begin{equation}
    \lambda_{\text{max}} = \max |\text{Eig}(T_{ij})|.
\end{equation}

\section{Recursive Feedback and Stochasticity}
Recursive feedback is defined as:
\begin{equation}
    f(t) = \alpha R(t) + \beta M(t),
\end{equation}
with stochastic dynamics incorporated as:
\begin{equation}
    \rho_k(t) \to \rho_k(t) + \xi_k(t),
\end{equation}
where $\xi_k(t)$ follows a Gaussian distribution.

\section{Educational and Visualization Models}
Fibonacci spirals are represented in 2D as:
\begin{equation}
    r_n = \frac{1}{\phi^n}, \quad \theta_n = 2\pi n,
\end{equation}
where $\phi = \frac{1 + \sqrt{5}}{2}$ is the golden ratio. Cartesian coordinates:
\begin{equation}
    x_n = r_n \cos(\theta_n), \quad y_n = r_n \sin(\theta_n).
\end{equation}
For 3D extensions:
\begin{equation}
    z_n = c_n \sin(\phi \cdot n).
\end{equation}

\section{Cryptographic Framework}
Entropy of prime-based keys:
\begin{equation}
    H = - \sum_{i} p_i \log_2 p_i.
\end{equation}
Quantum-resistant encryption:
\begin{equation}
    \text{Key}(t) = H(P(x, t)) \cdot Q_{\text{resistant}}.
\end{equation}

\section{Interdisciplinary Extensions}
\subsection{Astrophysics}
Galactic spiral structures:
\begin{equation}
    r_n = \frac{1}{\phi^n}, \quad \theta_n = 2\pi n,
\end{equation}
with corrections from observational data:
\begin{equation}
    r_{\text{obs}} = r_n \cdot (1 + \delta(t)).
\end{equation}
\subsection{Biology}
Phyllotaxis modeled using Fibonacci spirals:
\begin{equation}
    r_n = k \cdot n^{1/2}, \quad \theta_n = 2\pi n \cdot \phi.
\end{equation}
\section{Test Results Overview}

In this section, we summarize the results of the tests conducted on the properties and relationships between Pythagorean triplets and Fibonacci numbers. These tests were designed to validate the theoretical findings and provide empirical evidence supporting our hypotheses.

\subsection{Methodology}
The tests were conducted on several sets of data, including both small and large Pythagorean triplets and corresponding Fibonacci sequences. The primary metrics measured included:

\begin{itemize}
    \item The consistency of the Fibonacci sequence's appearance in the sides of Pythagorean triplets.
    \item The efficiency of algorithms for generating Pythagorean triplets from Fibonacci numbers.
    \item The performance of algorithms in terms of time complexity when generating large triplets and Fibonacci numbers.
\end{itemize}

\subsection{Key Findings}
\begin{itemize}
    \item The relationship between Fibonacci numbers and Pythagorean triplets was consistently observed, with Fibonacci triplets satisfying the Pythagorean theorem.
    \item Algorithms for generating Pythagorean triplets from Fibonacci sequences demonstrated a high level of efficiency, even for large input sizes.
    \item In certain cases, discrepancies were observed in the generation of large Fibonacci numbers, suggesting a need for optimization in the algorithm for better handling of larger datasets.
    \item The tests confirmed that certain Pythagorean triplets can be expressed as sums of squares of Fibonacci numbers, reinforcing the theoretical basis of the relationship.
\end{itemize}

\subsection{Summary}
Overall, the results of these tests validate the theoretical connections between Pythagorean triplets and Fibonacci numbers, offering empirical confirmation of their inherent link. These findings pave the way for further research into more efficient algorithms for generating and manipulating both sequences.

\section{Conclusion}
This enhanced framework integrates Pythagorean triplets, Fibonacci sequences, and Multiplicity Theory into a robust mathematical model. It demonstrates significant potential in cryptography, astrophysics, education, and interdisciplinary research, bridging theory with practical applications.

\nocite{*}
\bibliographystyle{unsrt}
\bibliography{references}

\end{document}